\documentclass[11pt]{article}

\usepackage[margin=1in]{geometry}
\usepackage{amsmath, amssymb, amsthm}
\usepackage{mathtools}
\usepackage[colorlinks=true, linkcolor=blue, urlcolor=blue, citecolor=blue]{hyperref}

\title{\textbf{Reading off a surface-code distance}}
\author{Sanjay Suresh}
\date{May 2026}

\begin{document}
\maketitle

The code distance $d$ of a planar surface code is just the length of the shortest logical
operator --- the shortest chain of data qubits connecting the two like-type boundaries. Eyeballing
it on the lattice is usually faster than any formula.

A distance-$d$ code corrects up to $\left\lfloor (d-1)/2 \right\rfloor$ arbitrary errors, so distance
grows the protection roughly exponentially \emph{below threshold}: doubling $d$ squares the
suppression of logical error. The catch is cost --- a planar patch needs on the order of $d^2$
physical qubits, so every step up in distance is quadratically expensive.

Practical takeaway: pick the smallest $d$ whose projected logical error rate clears your algorithm's
total operation budget, not the largest one you can fit.

\end{document}
